\chapter{Pleural Effusion Symptoms}

\section*{Definition}
Pleural effusion is the pathologic accumulation of fluid in the pleural space, classified as transudate or exudate by Light's criteria. Clinically significant effusions are those occupying more than 300 mL or causing respiratory compromise; up to 500 mL may be asymptomatic.

\section*{Pathophysiology}
Fluid accumulates when the rate of pleural fluid production (from systemic/hilar capillaries or across the pleura) exceeds lymphatic absorption. Transudative effusions result from increased hydrostatic pressure or decreased oncotic pressure, while exudative effusions arise from increased capillary permeability or impaired lymphatic drainage.

\section*{Etiology}
\begin{itemize}
  \item \textbf{Transudate:} Congestive heart failure (most common), cirrhosis with ascites, nephrotic syndrome, hypoalbuminemia, peritoneal dialysis, constrictive pericarditis
  \item \textbf{Exudate:} Pneumonia (parapneumonic empyema), malignancy (lung cancer, breast cancer, lymphoma), pulmonary embolism, tuberculosis, connective tissue disease (SLE, rheumatoid arthritis), pancreatitis, post-CABG
  \item \textbf{Hemothorax:} Trauma, aortic dissection, malignancy, anticoagulation
  \item \textbf{Chylothorax:} Thoracic duct injury, lymphoma, sarcoidosis, filariasis
\end{itemize}

\section*{Indications for Evaluation}
Seek care if:
\begin{itemize}
  \item Progressive dyspnea at rest or on minimal exertion
  \item Pleuritic chest pain or dull, heavy sensation on the affected side
  \item Fever, chills, or productive cough suggesting empyema
  \item Known malignancy with new-onset dyspnea
  \item Rapid accumulation of fluid (more than 1 L in 24 hours) causing mediastinal shift
\end{itemize}

\section*{Related Symptoms}
\begin{itemize}
  \item Dyspnea, pleuritic chest pain, decreased breath sounds, dullness to percussion, cough, orthopnea, fever
\end{itemize}
\textbf{Note:} The patient typically lies on the affected side to splint the chest and maximize ventilation of the unaffected lung.

\section*{Self-Care / Home Management}
\begin{itemize}
  \item Semi-upright positioning to optimize diaphragmatic excursion and ease breathing
  \item Oxygen supplementation if home oxygen is available and prescribed
  \item Monitor daily weight and report rapid changes to clinician
  \item Avoid heavy lifting or Valsalva maneuvers that may increase pleural pressure
  \item For heart failure patients: dietary sodium restriction (less than 2 g/day) and daily weight monitoring
\end{itemize}

\section*{Diagnostic Approach}
\begin{itemize}
  \item Chest X-ray (upright PA and lateral): blunting of costophrenic angle (200--500 mL), meniscus sign, opacification
  \item Ultrasound for detection of loculation, septations, and guidance for thoracentesis
  \item Diagnostic thoracentesis with Light's criteria: exudate if pleural/serum protein ratio greater than 0.5, LDH ratio greater than 0.6, or pleural LDH greater than 2/3 upper normal serum LDH
  \item Pleural fluid analysis: cell count/differential, Gram stain/culture, cytology, pH, glucose, triglycerides
  \item CT chest with contrast for suspected empyema, malignancy, or complex loculated effusion
\end{itemize}

\section*{Treatment / Management Overview}
\begin{itemize}
  \item Treat underlying cause (diuresis for CHF, antibiotics for pneumonia, anticoagulation for PE)
  \item Therapeutic thoracentesis for symptomatic large effusions; drain less than 1 L at a time to avoid re-expansion pulmonary edema
  \item Tube thoracostomy for empyema, hemothorax, or recurrent rapid reaccumulation
  \item Intrapleural fibrinolytics (tPA/DNase) for loculated parapneumonic effusions or empyema
  \item Pleurodesis (talc, doxycycline) or indwelling pleural catheter for recurrent malignant effusions
\end{itemize}

\clearpage
